\documentclass[aspectratio=169]{beamer}

% ---------- Font & ngôn ngữ (XeLaTeX, hỗ trợ tiếng Việt) ----------
\usepackage{fontspec}
\setmainfont{DejaVu Sans}
\setmonofont{DejaVu Sans Mono}
\usepackage{polyglossia}
\setdefaultlanguage{vietnamese}

\usepackage{amsmath,amssymb}
\usepackage{tikz}
\usetikzlibrary{arrows.meta,positioning,shapes.geometric,fit,backgrounds}
\usepackage{booktabs}
\usepackage{array}
\usepackage{graphicx}
\usepackage{multicol}

% ---------- Theme ----------
\usetheme{Madrid}
\usecolortheme{whale}
\setbeamertemplate{navigation symbols}{}
\setbeamertemplate{footline}[frame number]
\setbeamertemplate{caption}[numbered]

\title[Truck--Drone VRPTW]{Hệ thống lấy hàng truck và drone hoạt động song song, có ràng buộc time window}
\subtitle{Mô hình bài toán và giải thuật Tabu Search thích nghi}
\author{Học viên: Trần Xuân Nam -- 20251177M\\[1mm]
\small Giảng viên hướng dẫn: TS. Nguyễn Khánh Phương}
\date{}

% ---------- Tiện ích ----------
\newcommand{\hl}[1]{\textcolor{blue!60!black}{\textbf{#1}}}
% Khối pseudocode: chữ nhỏ, đơn cách
\newenvironment{pseudo}{\par\scriptsize\ttfamily\obeylines\obeyspaces}{\par}
\newcommand{\ind}{\hspace*{1.2em}}

\begin{document}

% ===================================================================
\begin{frame}
  \titlepage
\end{frame}

\begin{frame}{Nội dung}
  \begin{multicols}{2}
    \tableofcontents
  \end{multicols}
\end{frame}

% ===================================================================
\section{Giới thiệu bài toán}

\begin{frame}{Bối cảnh}
  \begin{columns}[T]
    \begin{column}{0.55\textwidth}
      \begin{itemize}
        \item Một \hl{Depot} xuất phát $K$ truck và $D$ drone giống nhau.
        \item Tập khách hàng $C$ cần được \hl{lấy hàng} (pickup), chia làm 2 loại:
        \begin{itemize}
          \item $C_1$: bắt buộc phục vụ bằng \textbf{truck}
          \item $C_2$: phục vụ được bằng \textbf{truck hoặc drone}
        \end{itemize}
        \item Mỗi xe được phép thực hiện \hl{nhiều chuyến (trip)} nối tiếp: xuất phát rỗng từ depot $\to$ lấy hàng $\to$ quay về depot.
      \end{itemize}
    \end{column}
    \begin{column}{0.45\textwidth}
      \begin{tikzpicture}[scale=0.9, every node/.style={font=\small}]
        \node[circle,draw,fill=orange!30,minimum size=8mm] (D) at (0,0) {Depot};
        \node[draw,fill=blue!20,rounded corners] (T) at (-2,-1.8) {Truck};
        \node[draw,fill=green!20,rounded corners] (Dr) at (2,-1.8) {Drone};
        \node[circle,draw,minimum size=6mm] (c1) at (-3.2,-3.2) {$i_1{\in}C_1$};
        \node[circle,draw,minimum size=6mm] (c2) at (-1.2,-3.2) {$i_2{\in}C_2$};
        \node[circle,draw,minimum size=6mm] (c3) at (1.2,-3.2) {$i_3{\in}C_2$};
        \node[circle,draw,minimum size=6mm] (c4) at (3.2,-3.2) {$i_4{\in}C_2$};
        \draw[-{Latex}] (D) -- (T);
        \draw[-{Latex}] (D) -- (Dr);
        \draw[-{Latex}] (T) -- (c1);
        \draw[-{Latex}] (T) -- (c2);
        \draw[-{Latex}] (Dr) -- (c3);
        \draw[-{Latex}] (Dr) -- (c4);
      \end{tikzpicture}
    \end{column}
  \end{columns}
\end{frame}

\begin{frame}{Ràng buộc chính và mục tiêu}
  \begin{itemize}
    \item \hl{Tải trọng}: truck $\le M_T$, drone $\le M_D$
    \item \hl{Tầm bay drone} $L_D$: tổng \emph{quãng đường} bay của \emph{một chuyến} drone (từ depot, đi lấy hàng, về depot) $\le L_D$
    \item \hl{Time window} khách $i$: $[e_i, l_i]$ — đến sớm phải chờ đến $e_i$, không được đến sau $l_i$
    \item \hl{Thời gian chờ hàng tối đa} $L_w$: thời điểm hàng về depot $-$ thời điểm lấy hàng $\le L_w$
    \item Thời gian sạc/thay pin drone và phục vụ khách $= 0$; drone luôn xuất phát với pin đầy
  \end{itemize}
  \vspace{2mm}
  \begin{block}{Mục tiêu}
    Xây dựng hành trình cho tất cả truck/drone sao cho \hl{tối thiểu hoá makespan} — thời điểm toàn bộ hàng hoá đã được đưa về depot (hệ thống bắt đầu tại thời điểm 0).
  \end{block}
\end{frame}

\begin{frame}{Ký hiệu đầu vào}
  \begin{columns}[T]
    \begin{column}{0.5\textwidth}
      \textbf{Khách hàng $i$}
      \begin{itemize}
        \item \small $q_i$: khối lượng hàng cần lấy
        \item \small $[e_i, l_i]$: cửa sổ thời gian
        \item \small $i \in C_1$: chỉ phục vụ bằng truck
        \item \small $i \in C_2$: truck hoặc drone
      \end{itemize}
      \textbf{Khoảng cách, thời gian}
      \begin{itemize}
        \item \small $d_{ij}$: khoảng cách Euclid từ $i$ đến $j$
        \item \small $\tau^v_{ij} = d_{ij} / \text{vận tốc}(v)$: thời gian đi bằng phương tiện $v$
      \end{itemize}
    \end{column}
    \begin{column}{0.5\textwidth}
      \textbf{Phương tiện}
      \begin{itemize}
        \item \small $K$, $D$: số truck, số drone
        \item \small $M_T$, $M_D$: tải trọng tối đa
        \item \small $L_D$: quãng đường tối đa một chuyến drone
        \item \small $L_w$: thời gian tối đa từ lúc lấy hàng đến lúc hàng về depot
      \end{itemize}
      \begin{alertblock}{Quy đổi dữ liệu test}
        \small Bộ dữ liệu test cho \texttt{drone\_lim} là \emph{thời gian bay} tối đa. Khi đọc dữ liệu:
        $$L_D = \texttt{drone\_lim} \times v_D$$
        Ví dụ: $700 \times 31.29 \approx 21\,905$.
      \end{alertblock}
    \end{column}
  \end{columns}
\end{frame}

% ===================================================================
\section{Biểu diễn nghiệm}

\begin{frame}{Cấu trúc dữ liệu nghiệm}
  \begin{columns}[T]
    \begin{column}{0.5\textwidth}
      \textbf{Một chuyến (trip):}
      $$\sigma_{vk} = [0, i_1, i_2, \ldots, i_m, 0]$$
      \textbf{Chuỗi chuyến của xe $v$:}
      $$R_v = (\sigma_{v1}, \sigma_{v2}, \ldots, \sigma_{vm_v})$$
      \begin{alertblock}{Lưu ý}
        Thứ tự chuyến là một phần của nghiệm — đổi thứ tự làm thay đổi thời gian bắt đầu/kết thúc của \emph{mọi} chuyến phía sau.
      \end{alertblock}
    \end{column}
    \begin{column}{0.5\textwidth}
      \small
      \textbf{Customer}: id, demand, $[e_i,l_i]$, droneEligible

      \textbf{Trip}: uid, vehicleId, customers[], startTime, returnTime, load, travelDistance, arrivalTime[], waitingTime[]

      \textbf{Vehicle}: id, type (TRUCK/DRONE), trips[], completionTime

      \textbf{Solution}: vehicles[], makespan, normalizedMakespan, $V_Q,V_D,V_{TW},V_W$, totalViolation, penalizedObjective, totalDistance
    \end{column}
  \end{columns}
\end{frame}

\begin{frame}{Tính lịch cho một phương tiện (RECOMPUTE\_VEHICLE)}
  \small Hàm nền tảng: mọi toán tử đều gọi hàm này sau khi thay đổi nghiệm.

  \textbf{Thời gian bắt đầu chuyến:}
  $$t^{\text{start}}_{v,1} = 0, \qquad t^{\text{start}}_{v,k} = r_{v,k-1}$$
  \textbf{Thời gian lấy hàng} với $\sigma = [0, i_1, \ldots, i_m, 0]$:
  $$a_{i_1} = \max\{e_{i_1},\, t^{\text{start}}_\sigma + \tau^v_{0 i_1}\}, \qquad
    a_{i_h} = \max\{e_{i_h},\, a_{i_{h-1}} + \tau^v_{i_{h-1} i_h}\}$$
  Nếu $a_i > l_i$: vẫn lấy hàng tại $a_i$, nhưng phát sinh vi phạm $a_i - l_i$.

  \textbf{Thời gian trở về depot và thời gian chờ hàng:}
  $$r_\sigma = a_{i_m} + \tau^v_{i_m 0}, \qquad w_i = r_\sigma - a_i, \qquad C_v = r_{\text{chuyến cuối}}$$

  \begin{center}
  \begin{tikzpicture}[node distance=6mm, every node/.style={font=\footnotesize}]
    \node[draw,rounded corners,fill=orange!20] (d0) {Depot: $t^{\text{start}}$};
    \node[draw,rounded corners,fill=blue!15,right=of d0] (i1) {$i_1$: $a_{i_1}$};
    \node[draw,rounded corners,fill=blue!15,right=of i1] (i2) {$i_2$: $a_{i_2}$};
    \node[draw,rounded corners,fill=blue!15,right=of i2] (i3) {$\cdots$};
    \node[draw,rounded corners,fill=orange!20,right=of i3] (d1) {Depot: $r_\sigma$};
    \draw[-{Latex}] (d0)--(i1); \draw[-{Latex}] (i1)--(i2);
    \draw[-{Latex}] (i2)--(i3); \draw[-{Latex}] (i3)--(d1);
  \end{tikzpicture}
  \end{center}
\end{frame}

\begin{frame}{Kiểm tra tương thích tĩnh (STATIC\_COMPATIBLE)}
  \small Mỗi cặp (khách $i$, xe $v$) được lọc bởi một điều kiện \hl{tĩnh}, không phụ thuộc nghiệm hiện tại — dùng ở mọi nơi sinh move và trong Ruin \& Recreate.

  \begin{columns}[T]
    \begin{column}{0.5\textwidth}
      \textbf{Với xe DRONE:}
      \begin{itemize}
        \item \small $i \in C_2$
        \item \small $q_i \le M_D$
        \item \small $d_{0i} + d_{i0} \le L_D$ (quãng đường khứ hồi)
        \item \small $\tau^{D}_{i0} \le L_w$
      \end{itemize}
    \end{column}
    \begin{column}{0.5\textwidth}
      \textbf{Với xe TRUCK:}
      \begin{itemize}
        \item \small $q_i \le M_T$
        \item \small $\tau^{T}_{i0} \le L_w$
      \end{itemize}
    \end{column}
  \end{columns}

  \vspace{2mm}
  \begin{alertblock}{Ý nghĩa}
    \small Khách $i$ phục vụ khả thi được bằng $v$ ít nhất trong chuyến đơn $[0,i,0]$. Tải trọng, tầm bay, thời gian chờ được phép vi phạm tạm thời ở cấp chuyến, nhưng không đưa khách vào loại xe mà nó \emph{không bao giờ} có thể được phục vụ khả thi.
  \end{alertblock}
\end{frame}

% ===================================================================
\section{Hàm đánh giá}

\begin{frame}{Bốn loại vi phạm chuẩn hoá}
  Ký hiệu $[x]^+ = \max\{0,x\}$ và hằng số thời gian $H = \max\{1, C_{\max}(s_0), \max_i l_i\}$ — tính \hl{một lần} từ nghiệm ban đầu, không đổi.

  \begin{align*}
    V_Q(s)  &= \frac{1}{n}\sum_{\sigma} \frac{[Q(\sigma) - M_{v(\sigma)}]^+}{M_{v(\sigma)}}
      & \text{\small vượt \hl{tải trọng}}, \ Q(\sigma)=\textstyle\sum_{i\in\sigma} q_i\\[1mm]
    V_D(s)  &= \frac{1}{n}\sum_{\sigma \in \mathcal{R}_D(s)} \frac{[L(\sigma) - L_D]^+}{L_D}
      & \text{\small vượt \hl{tầm bay drone}}\\[1mm]
    V_{TW}(s) &= \frac{1}{n}\sum_{i \in C} \frac{[a_i - l_i]^+}{H}
      & \text{\small trễ \hl{time window}}\\[1mm]
    V_W(s)  &= \frac{1}{n}\sum_{i \in C} \frac{[r_{\sigma(i)} - a_i - L_w]^+}{L_w}
      & \text{\small vượt \hl{thời gian chờ hàng}}
  \end{align*}

  \small $L(\sigma)$: tổng quãng đường của chuyến drone $\sigma$, tính cả đoạn từ depot và quay về depot.
\end{frame}

\begin{frame}{Tổng vi phạm và điều kiện khả thi}
  $$V_\Sigma(s) = V_Q(s) + V_D(s) + V_{TW}(s) + V_W(s)$$

  \begin{block}{Nghiệm khả thi}
    $$V_\Sigma(s) \le \varepsilon \quad \text{\textbf{và}} \quad \text{unassigned}(s) = 0, \qquad \varepsilon = 10^{-9}$$
  \end{block}

  \begin{alertblock}{Nguyên tắc cốt lõi của thuật toán}
    Trong quá trình tìm kiếm, nghiệm được phép \hl{tạm thời vi phạm} cả 4 loại ràng buộc trên. \\
    Nghiệm \textbf{đầu ra} bắt buộc thoả tất cả ràng buộc.
  \end{alertblock}

  \small Ràng buộc \emph{cấu trúc} không bao giờ được vi phạm: khách không trùng/thiếu, $C_1$ không nằm trên drone, trip phải bắt đầu/kết thúc ở depot, khách phải tương thích tĩnh với phương tiện.

  \vspace{1mm}
  \small \textbf{Lưu ý triển khai:} cần đếm riêng số khách chưa gán (\texttt{unassignedCount}) và bắt buộc bằng 0; $V_\Sigma(s)=0$ không tự đảm bảo điều này.
\end{frame}

\begin{frame}{Makespan chuẩn hoá và hàm mục tiêu có phạt}
  \textbf{Makespan:}
  $$C_{\max}(s) = \max_v C_v(s), \qquad \bar{C}_{\max}(s) = \frac{C_{\max}(s)}{H}$$

  \textbf{Hàm đánh giá có phạt} (chỉ dùng để \emph{dẫn dắt} tìm kiếm):
  $$F_\lambda(s) = \bar{C}_{\max}(s) + \lambda_Q V_Q(s) + \lambda_D V_D(s) + \lambda_{TW} V_{TW}(s) + \lambda_W V_W(s)$$

  \begin{alertblock}{Giá trị mục tiêu thật sự}
    Mục tiêu cuối cùng luôn là \hl{makespan thực tế $C_{\max}$} — $F_\lambda$ chỉ hỗ trợ định hướng vì $\lambda$ thay đổi liên tục trong quá trình tìm kiếm (Strategic Oscillation).
  \end{alertblock}

  \small Chia cho $n$ và $H$ giúp các thành phần cùng thang đo, $\lambda$ không phụ thuộc cỡ bài.
\end{frame}

% ===================================================================
\section{Khởi tạo nghiệm}

\begin{frame}{Xây dựng nghiệm khởi tạo (Greedy Insertion)}
  \textbf{Ý tưởng:} chèn tuần tự khách theo deadline tăng dần; mọi xe ban đầu chưa có chuyến nào.

  \begin{enumerate}
    \item Sắp khách hàng theo $l_i$ tăng dần
    \item Với mỗi khách $i$, trên mọi xe tương thích, sinh mọi move:
    \begin{itemize}
      \item \small mọi vị trí trong mọi trip đã có của xe
      \item \small một trip mới $[0,i,0]$ tại mọi vị trí trong chuỗi trip của xe
    \end{itemize}
    \item Tạm chèn, tính lại lịch trip bị ảnh hưởng và các trip sau; loại move vi phạm ràng buộc cứng
    \item Phân loại:
    \begin{itemize}
      \item \small \textbf{feasibleMoves} (không tạo thêm vi phạm TW) $\to$ chọn makespan nhỏ nhất
      \item \small \textbf{penalizedMoves} $\to$ chọn theo thứ tự từ điển: (1) $\Delta TW$ nhỏ nhất, (2) makespan nhỏ nhất, (3) quãng đường tăng ít nhất
    \end{itemize}
    \item Chèn vĩnh viễn move được chọn; sau cùng tính $H$
  \end{enumerate}
\end{frame}

\begin{frame}[fragile]{Hàm $\Delta TW(a)$}
  \begin{columns}[T]
    \begin{column}{0.48\textwidth}
      \begin{block}{compute $\Delta TW(a)$}
        \begin{pseudo}
1. penalty $\leftarrow$ 0
2. for mỗi khách j đã được chèn
\ind trong solution x của move a do
3. \ind penalty $\leftarrow$ penalty + max(0, t\_j - l\_j)
4. end for
5. return penalty
        \end{pseudo}
      \end{block}
      \small $\Delta TW(a)$ = phần trễ \hl{tăng thêm} so với trước khi chèn.
    \end{column}
    \begin{column}{0.52\textwidth}
      \textbf{Vì sao tính trên mọi khách?}
      \begin{itemize}
        \item \small Chèn $i$ vào giữa chuyến làm \hl{mọi khách sau $i$} đến muộn hơn.
        \item \small Mọi \hl{chuyến sau} của cùng xe cũng bắt đầu muộn hơn.
        \item \small Nếu chỉ xét riêng $i$, một move đẩy khách khác ra khỏi time window vẫn bị xếp nhầm vào feasibleMoves.
      \end{itemize}
      \textbf{Cài đặt:} chỉ xe đích thay đổi, nên
      $$\Delta TW = \sum_{j \in v} [a_j - l_j]^+_{\text{sau}} - \sum_{j \in v} [a_j - l_j]^+_{\text{trước}}$$
    \end{column}
  \end{columns}
\end{frame}

% ===================================================================
\section{Tabu Search tổng quan}

\begin{frame}{Sơ đồ tổng quát}
  \begin{center}
  \resizebox{\linewidth}{!}{%
  \begin{tikzpicture}[
      node distance=4mm and 6mm,
      box/.style={draw, rounded corners, fill=blue!12, align=center, font=\scriptsize, minimum height=9mm, text width=2.1cm},
      dbox/.style={draw, rounded corners, dashed, fill=orange!12, align=center, font=\scriptsize, minimum height=9mm, text width=2.1cm},
      dia/.style={draw, diamond, fill=green!12, align=center, font=\scriptsize, inner sep=1pt, text width=1.7cm},
      term/.style={draw, rounded corners, fill=gray!15, align=center, font=\scriptsize, minimum height=9mm, text width=2.1cm}
    ]
    \node[box] (init) {Nghiệm khởi tạo \& TS init};
    \node[box,right=of init] (gen) {Sinh \& đánh giá neighborhood moves};
    \node[box,right=of gen] (sel) {Chọn move tốt nhất chấp nhận được};
    \node[box,right=of sel] (upd) {Cập nhật nghiệm, tabu, penalty, incumbent};
    \node[dia,right=of upd] (stop) {Dừng?};
    \node[term,right=of stop] (ret) {Trả về $s^*$};
    \node[dia,below=of upd] (stag) {Trì trệ?\\$h_{Div}{\ge}H_{Div}$};
    \node[dbox,below=of sel] (rr) {Diversification\\Ruin-\&-Recreate};

    \draw[-{Latex}] (init) -- (gen);
    \draw[-{Latex}] (gen) -- (sel);
    \draw[-{Latex}] (sel) -- (upd);
    \draw[-{Latex}] (upd) -- (stop);
    \draw[-{Latex}] (stop) -- node[above,font=\scriptsize]{Yes} (ret);
    \draw[-{Latex}] (stop) -- node[right,font=\scriptsize]{No} (stag);
    \draw[-{Latex}] (stag) -- node[above,font=\scriptsize]{Yes} (rr);
    \draw[-{Latex}] (rr) -- (gen);
    \draw[-{Latex}] (stag.west) to[out=180,in=-90] node[pos=0.35,left,font=\scriptsize]{No} (gen.south);
  \end{tikzpicture}%
  }
  \end{center}

  \small
  \hl{Vòng lặp chính}: sinh candidate pool $\to$ lọc theo tabu/aspiration $\to$ chọn move tốt nhất (có thể xấu hơn hiện tại!) $\to$ cập nhật nghiệm, tabu, $\lambda$, nghiệm tốt nhất.
\end{frame}

\begin{frame}{Các nghiệm cần lưu}
  \small Thuật toán duy trì đồng thời \hl{ba loại nghiệm}.

  \vspace{2mm}
  \textbf{Nghiệm hiện tại $s$}\\
  \small Nghiệm mà Tabu Search đang đứng tại đó. $s$ có thể khả thi hoặc không khả thi.

  \vspace{2mm}
  \textbf{Nghiệm khả thi tốt nhất $s^*$}\\
  \small Chỉ cập nhật khi $V_\Sigma(s) = 0$, \texttt{unassigned}$(s)=0$ \emph{và} $C_{\max}(s) < C_{\max}(s^*)$. Nếu chưa tìm được nghiệm khả thi nào thì \texttt{bestFeasible} $\leftarrow$ \texttt{null}.

  \vspace{2mm}
  \textbf{Nghiệm không khả thi tốt nhất $s^I$}\\
  \small Hai nghiệm không khả thi được so sánh theo thứ tự từ điển (BETTER\_INFEASIBLE):
  $$\big(V_\Sigma(s), \bar{C}_{\max}(s), \text{Distance}(s)\big)$$
  \small Ưu tiên: (1) tổng vi phạm nhỏ hơn, (2) nếu bằng nhau thì makespan nhỏ hơn, (3) nếu vẫn bằng nhau thì tổng quãng đường nhỏ hơn.
\end{frame}

\begin{frame}{Cách sử dụng ba loại nghiệm}
  \begin{itemize}
    \item $s$ là điểm xuất phát để \hl{sinh candidate pool} ở mỗi vòng lặp — luôn di chuyển, kể cả khi move được chọn xấu hơn $s$ hiện tại.
    \item $s^*$ là \hl{kết quả trả về} và \hl{tiêu chuẩn aspiration} khi đã có nghiệm khả thi: chỉ chấp nhận move tabu nếu candidate khả thi và cải thiện $C_{\max}(s^*)$.
    \item $s^I$ định hướng khi \hl{chưa có nghiệm khả thi nào} — cho biết đang tiến gần vùng khả thi hay không, và là chuẩn so sánh cho aspiration ở giai đoạn này.
    \item Sau mỗi vòng lặp, \texttt{UPDATE\_BEST\_SOLUTIONS} cập nhật $s^*$, $s^I$ và bộ đếm $h_{Div}$, $h_{Stop}$.
    \item Khi dừng: $s^* \ne$ null $\to$ trả về $s^*$; ngược lại báo ``No feasible solution found'' và trả $s^I$ chỉ mang tính chẩn đoán.
  \end{itemize}
\end{frame}

% ===================================================================
\section{Sinh lân cận}

\begin{frame}{Chọn thành phần tìm kiếm (SELECT\_SEARCH\_COMPONENTS)}
  \small Phiên bản đầu sinh \hl{toàn bộ move} để kiểm tra tính đúng đắn; sau đó giới hạn nguồn sinh move để giảm thời gian tính.

  \begin{columns}[T]
    \begin{column}{0.5\textwidth}
      \textbf{Khi chưa có nghiệm khả thi}
      \begin{itemize}
        \item \small Khách đóng góp lớn vào $V_{TW}$, $V_W$
        \item \small Khách thuộc trip vượt tải
        \item \small Khách thuộc trip drone vượt tầm bay
        \item \small Một tỷ lệ khách ngẫu nhiên (30\%)
        \item \small Trip = các trip chứa những khách trên
      \end{itemize}
    \end{column}
    \begin{column}{0.5\textwidth}
      \textbf{Khi đã có nghiệm khả thi}
      $$\mathcal{V}^{\text{crit}} = \{v : C_v(s) = C_{\max}(s)\}$$
      \begin{itemize}
        \item \small Khách và trip trên xe critical
        \item \small Thêm 30\% khách và 30\% trip ngẫu nhiên toàn hệ thống
      \end{itemize}
    \end{column}
  \end{columns}

  \vspace{2mm}
  \begin{block}{Lưu ý}
    \small Giới hạn chỉ áp dụng cho \emph{nguồn} của move. \hl{Đích} của move luôn được xét trên tất cả phương tiện và trip tương thích.
  \end{block}
\end{frame}

\begin{frame}[fragile]{Sinh candidate pool (BUILD\_CANDIDATE\_POOL)}
  \begin{columns}[T]
    \begin{column}{0.5\textwidth}
      \begin{block}{Pseudocode}
        \begin{pseudo}
(cust, trips) $\leftarrow$ SELECT\_SEARCH\_COMPONENTS(s, s*)
rawMoves $\leftarrow$ GENERATE\_RELOCATE(s, cust)
\ind $\cup$ GENERATE\_OR\_OPT\_2(s, trips)
\ind $\cup$ GENERATE\_SWAP(s, cust)
\ind $\cup$ GENERATE\_2\_OPT(s, trips)
\ind $\cup$ GENERATE\_CROSS\_TRIP(s, trips)
\ind $\cup$ GENERATE\_TRIP\_RELOCATE(s, trips)
for each move m in rawMoves:
\ind c $\leftarrow$ EVALUATE\_MOVE(s, m, $\lambda$)
\ind if c is valid: add c to pool
return pool
        \end{pseudo}
      \end{block}
    \end{column}
    \begin{column}{0.5\textwidth}
      \textbf{Giới hạn số move khi instance lớn} (lấy mẫu ngẫu nhiên; instance nhỏ vẫn duyệt đủ):
      \begin{center}
      \small
      \begin{tabular}{ll}
        \toprule
        Toán tử & Tối đa \\
        \midrule
        Relocate & 80 đích / khách \\
        Or-opt(2) & 80 đích / block \\
        Swap & 600 cặp \\
        Cross-trip & 40 / cặp trip, 800 / vòng \\
        2-opt, Trip-relocate & duyệt đủ \\
        \bottomrule
      \end{tabular}
      \end{center}
    \end{column}
  \end{columns}
\end{frame}

\begin{frame}{1) Relocate \quad 2) Or-opt(2)}
  \begin{columns}[T]
    \begin{column}{0.5\textwidth}
      \textbf{Relocate}
      \begin{itemize}
        \item \small Lấy 1 khách $i$ khỏi vị trí hiện tại $[\ldots,p,i,q,\ldots]$.
        \item \small Chèn vào mọi vị trí trong mọi trip tương thích, hoặc một trip mới tại mọi vị trí trong chuỗi trip.
        \item \small Xoá cung $(p,i),(i,q)$, sinh cung $(p,q)$ tại nguồn; sinh cung mới tại vị trí chèn.
      \end{itemize}
      \small\textit{Ví dụ:} T1 $[0,3,5,8,0]$, D1 $[0,2,0]$ \\ $\to$ T1 $[0,3,8,0]$, D1 $[0,2,5,0]$
    \end{column}
    \begin{column}{0.5\textwidth}
      \textbf{Or-opt(2)}
      \begin{itemize}
        \item \small Di chuyển 1 block gồm 2 khách liên tiếp $[i_h, i_{h+1}]$ sang vị trí khác.
        \item \small Giữ nguyên thứ tự trong block khi chèn.
        \item \small Xe đích phải tương thích với \emph{cả hai} khách.
      \end{itemize}
      \small\textit{Ví dụ:} T1 $[0,3,5,8,0]$, T2 $[0,4,0]$ \\ $\to$ T1 $[0,8,0]$, T2 $[0,4,3,5,0]$
    \end{column}
  \end{columns}
\end{frame}

\begin{frame}{3) Swap \quad 4) 2-opt}
  \begin{columns}[T]
    \begin{column}{0.5\textwidth}
      \textbf{Swap}
      \begin{itemize}
        \item \small Hoán đổi 2 khách bất kỳ $i, j$ (có thể khác trip, khác xe).
        \item \small Điều kiện: $i$ tương thích tĩnh với xe của $j$ \emph{và} $j$ tương thích tĩnh với xe của $i$.
        \item \small Bỏ qua nếu hoán đổi không thay đổi gì.
      \end{itemize}
      \small\textit{Ví dụ:} T1 $[0,3,5,0]$, D1 $[0,2,0]$ \\ $\to$ T1 $[0,3,2,0]$, D1 $[0,5,0]$
    \end{column}
    \begin{column}{0.5\textwidth}
      \textbf{2-opt}\\
      \small Đảo ngược một đoạn trong cùng một trip:
      $$[\ldots,i_p,\ldots,i_q,\ldots] \to [\ldots,i_q,\ldots,i_p,\ldots]$$
      \begin{itemize}
        \item \small Chỉ trong nội bộ 1 trip; tập khách và tải không đổi.
        \item \small Gỡ các cung đi chéo, rút ngắn quãng đường và thời gian về depot.
      \end{itemize}
      \small\textit{Ví dụ:} $[0,3,5,8,6,0] \to [0,8,5,3,6,0]$
    \end{column}
  \end{columns}
\end{frame}

\begin{frame}{5) Cross-trip \quad 6) Trip-relocate}
  \begin{columns}[T]
    \begin{column}{0.5\textwidth}
      \textbf{Cross-trip}\\
      \small Chọn 2 trip, trao đổi 2 đoạn đuôi:
      \begin{align*}
        \sigma_a &= [A_1|A_2], & \sigma_b &= [B_1|B_2] \\
        \to\ \sigma_a' &= [A_1|B_2], & \sigma_b' &= [B_1|A_2]
      \end{align*}
      \small Có thể cùng xe, khác xe, hoặc giữa truck \& drone (nếu khách chuyển sang drone đều tương thích tĩnh).
    \end{column}
    \begin{column}{0.5\textwidth}
      \textbf{Trip-relocate}\\
      \small Di chuyển \emph{nguyên cả một trip} sang:
      \begin{itemize}
        \item \small vị trí khác trong chuỗi trip của cùng xe
        \item \small một xe tương thích khác, tại bất kỳ vị trí nào
      \end{itemize}
      \small Đặc thù bài toán nhiều chuyến: chuyển việc từ xe critical sang xe rảnh giảm $C_{\max}$ trực tiếp.
    \end{column}
  \end{columns}
  \vspace{3mm}
  \begin{alertblock}{Chung cho cả 6 toán tử}
    Mọi move được đánh giá tạm thời trên bản sao $\to$ tính lại lịch các xe bị ảnh hưởng $\to$ EVALUATE\_SOLUTION $\to$ loại nếu vi phạm ràng buộc \emph{cấu trúc}.
  \end{alertblock}
\end{frame}

\begin{frame}{Tạo trip mới \& huỷ trip rỗng}
  \small
  \textbf{Tạo trip mới}
  \begin{itemize}
    \item Relocate và Or-opt(2) cho phép \hl{tạo một trip mới} tại \emph{mọi vị trí} trong chuỗi trip $R_v$ của xe đích tương thích; Trip-relocate đặt trip được chuyển vào mọi vị trí của $R_v$.
    \item Vị trí trip trong $R_v$ là một phần của move: chèn ở vị trí $k$ làm dịch thời gian bắt đầu mọi trip từ $k{+}1$ trở đi.
    \item Ví dụ: $R_v = (\sigma_{v1}, \sigma_{v2}) \to R_v' = (\sigma_{v1}, [0,i,0], \sigma_{v2})$.
  \end{itemize}

  \vspace{1mm}
  \textbf{Huỷ trip rỗng}
  \begin{itemize}
    \item Sau khi áp dụng move, trip nguồn bị rút hết khách $\to$ \hl{xoá trip đó} khỏi $R_v$ trước khi tính lại lịch.
    \item Trip rỗng không mang thông tin nghiệm nhưng vẫn chiếm một vị trí trong $R_v$, làm các move sau sinh move vô nghĩa hoặc trùng lặp.
  \end{itemize}

  \begin{alertblock}{Ví dụ: Relocate khách $i$ ra khỏi $[\ldots,p,i,q,\ldots]$}
    \small Trip còn khách khác: mất cung $(p,i),(i,q)$, sinh cung $(p,q)$. \quad $i$ là khách duy nhất: $[0,i,0] \to$ rỗng $\to$ \hl{bị xoá}.
  \end{alertblock}
\end{frame}

% ===================================================================
\section{Đánh giá và chọn move}

\begin{frame}{Đánh giá một move (EVALUATE\_MOVE)}
  \small Một move phải được áp dụng \hl{tạm thời} (trên bản sao), sau đó tính lại lịch của các phương tiện bị ảnh hưởng.

  \vspace{2mm}
  \begin{enumerate}
    \item \small $s' \leftarrow$ bản sao của $s$; trích thuộc tính tabu cũ \texttt{oldAttributes}
    \item \small Áp dụng move lên $s'$; \hl{xoá mọi trip rỗng}
    \item \small Nếu $s'$ vi phạm ràng buộc cấu trúc $\to$ \texttt{INVALID}
    \item \small \texttt{RECOMPUTE\_VEHICLE} cho các xe bị ảnh hưởng
    \item \small \texttt{EVALUATE\_SOLUTION}$(s', \lambda)$; nếu vượt ngưỡng vi phạm tối đa cho phép $\to$ \texttt{INVALID}
    \item \small Trích \texttt{newAttributes}; trả về candidate với \\
      \texttt{removed} $=$ old $\setminus$ new, \quad \texttt{added} $=$ new $\setminus$ old
  \end{enumerate}

  \vspace{1mm}
  \begin{alertblock}{Ràng buộc cấu trúc phải loại ngay}
    \small Khách xuất hiện 0 hoặc $\ge 2$ lần; $C_1$ nằm trên drone; dùng phương tiện ngoài $K+D$; trip không bắt đầu/kết thúc tại depot; khách giao cho phương tiện không tương thích tĩnh. \\
    Vượt tải, vượt tầm bay, trễ TW, vượt $L_w$ \emph{không} bị loại — chỉ bị phạt.
  \end{alertblock}
\end{frame}

\begin{frame}{Danh sách Tabu — Thuộc tính}
  \textbf{Thuộc tính cung (Arc)}\\
  \small Mỗi cung: $(\text{vehicleId}, \text{fromNode}, \text{toNode})$ — tính cả cung Depot$\to$khách đầu, giữa hai khách, khách cuối$\to$Depot.\\
  Ví dụ: Relocate khách $i$ khỏi $[\ldots,p,i,q,\ldots]$ xoá $(p,i),(i,q)$ và tạo $(p,q)$. Move đảo ngược phải tái tạo $(p,i)$ và $(i,q)$ $\to$ bị tabu.

  \vspace{2mm}
  \textbf{Thuộc tính phân công (Assignment)}\\
  \small Khách $i$ chuyển từ xe $v_a$ sang $v_b$: lưu \texttt{("ASSIGN", i, }$v_a$\texttt{)} — trong thời hạn tabu, $i$ không được quay lại $v_a$, trừ khi thoả aspiration.

  \vspace{2mm}
  \textbf{Thuộc tính thứ tự trip (Trip Order)}\\
  \small Với Trip-relocate: lưu \texttt{("TRIP\_RETURN", tripUid, sourceVehicleId, sourcePredecessorTripUid)} — ngăn đưa trip quay ngay về vị trí cũ.
\end{frame}

\begin{frame}{Danh sách Tabu — Tenure \& kiểm tra}
  \textbf{Tabu tenure}
  $$\tau \sim U[\tau_0, 2\tau_0], \qquad \tau_0 = 7$$
  \small \texttt{tenure} $\leftarrow$ \texttt{RANDOM\_INTEGER(tau0, 2*tau0)}\\
  \small \texttt{tabuUntil[attribute]} $\leftarrow$ \texttt{currentIteration + tenure}

  \vspace{3mm}
  \textbf{Kiểm tra tabu (IS\_TABU)}
  \begin{itemize}
    \item \small Với mỗi thuộc tính $a$ trong \texttt{candidate.addedAttributes}: nếu \texttt{tabuUntil[a] > iteration} $\to$ candidate bị tabu.
  \end{itemize}

  \vspace{3mm}
  \textbf{Ghi tabu sau khi thực hiện move (REGISTER\_TABU)}
  \begin{itemize}
    \item \small Sinh tenure ngẫu nhiên mới, gán \texttt{tabuUntil[a]} $\leftarrow$ \texttt{iteration + tenure} cho mỗi thuộc tính vừa bị \hl{xoá} bởi move đã chọn (\texttt{removedAttributes}).
  \end{itemize}
\end{frame}

\begin{frame}{Aspiration criterion}
  \small Move tabu vẫn được phép thực hiện trong hai trường hợp.

  \vspace{2mm}
  \textbf{Đã có nghiệm khả thi}\\
  \small Candidate khả thi \emph{và} cải thiện makespan tốt nhất:
  $$V_\Sigma(s') = 0 \quad \text{và} \quad C_{\max}(s') < C_{\max}(s^*)$$

  \vspace{1mm}
  \textbf{Chưa có nghiệm khả thi}\\
  \small Candidate khả thi (nghiệm khả thi đầu tiên), hoặc cải thiện nghiệm không khả thi tốt nhất:
  $$\big(V_\Sigma(s'), C_{\max}(s'), \text{Distance}(s')\big) < \big(V_\Sigma(s^I), C_{\max}(s^I), \text{Distance}(s^I)\big)$$

  \vspace{1mm}
  \begin{alertblock}{Không dùng điều kiện $F_\lambda(s') < F_\lambda(s^*)$}
    \small Vì hệ số phạt $\lambda$ thay đổi liên tục (Strategic Oscillation), so sánh theo $F_\lambda$ giữa hai thời điểm không còn ý nghĩa nhất quán.
  \end{alertblock}
\end{frame}

\begin{frame}{Chọn move (SELECT\_BEST\_CANDIDATE)}
  \textbf{Khi chưa tìm được nghiệm khả thi}\\
  \small Mục tiêu chính: đưa nghiệm về vùng khả thi. Chọn move theo thứ tự từ điển:
  $$\big(V_\Sigma(s'), \bar{C}_{\max}(s'), \text{Distance}(s')\big)$$

  \vspace{2mm}
  \textbf{Khi đã có nghiệm khả thi}\\
  \small Chọn theo:
  $$\big(F_\lambda(s'), V_\Sigma(s'), \bar{C}_{\max}(s'), \text{Distance}(s')\big)$$

  \vspace{2mm}
  \begin{block}{Lưu ý}
    \small Chọn trong \texttt{admissiblePool} (không tabu, hoặc thoả aspiration). \hl{Move được chọn có thể xấu hơn nghiệm hiện tại} — đặc trưng giúp Tabu Search thoát khỏi tối ưu cục bộ.
  \end{block}
\end{frame}

% ===================================================================
\section{Thích nghi và đa dạng hoá}

\begin{frame}[fragile]{Cập nhật nghiệm tốt nhất và bộ đếm trì trệ}
  \begin{columns}[T]
    \begin{column}{0.55\textwidth}
      \begin{block}{UPDATE\_BEST\_SOLUTIONS}
        \begin{pseudo}
improved $\leftarrow$ false
if IS\_FEASIBLE(s):
\ind if s* = null or C(s) < C(s*) - $\varepsilon$:
\ind\ind s* $\leftarrow$ s; improved $\leftarrow$ true
else if BETTER\_INFEASIBLE(s, s$^I$):
\ind s$^I$ $\leftarrow$ s
\ind if s* = null: improved $\leftarrow$ true
return improved
        \end{pseudo}
      \end{block}
    \end{column}
    \begin{column}{0.45\textwidth}
      \textbf{Hai bộ đếm}
      \begin{itemize}
        \item \small $h_{Div}$: số vòng từ lần cải thiện hoặc đa dạng hoá gần nhất $\to$ kích hoạt Ruin \& Recreate khi $\ge H_{Div}$
        \item \small $h_{Stop}$: số vòng từ lần cải thiện record gần nhất $\to$ dừng khi $\ge H_{Stop}$
      \end{itemize}
      \textbf{Cập nhật}
      \begin{itemize}
        \item \small improved $\to h_{Div} \leftarrow 0,\ h_{Stop} \leftarrow 0$
        \item \small ngược lại $\to$ cả hai tăng 1
      \end{itemize}
    \end{column}
  \end{columns}
\end{frame}

\begin{frame}{Strategic Oscillation}
  \small Mỗi ràng buộc $g \in \{Q, D, TW, W\}$ có hệ số phạt riêng $\lambda_g$. Sau mỗi vòng lặp, nếu $V_g(s) \le \varepsilon$ thì \texttt{feasibleCounts[g] += 1}. Cứ mỗi $L_\lambda$ vòng:
  $$\rho_g = \frac{\#\{s_t : V_g(s_t) = 0\}}{L_\lambda}, \qquad
  \lambda_g \leftarrow
  \begin{cases}
    \min\{\lambda_g^{\max}, \gamma_\uparrow \lambda_g\} & \rho_g < \rho_{\min} \\
    \max\{\lambda_g^{\min}, \gamma_\downarrow \lambda_g\} & \rho_g > \rho_{\max} \\
    \lambda_g & \text{ngược lại}
  \end{cases}$$
  rồi đặt lại \texttt{feasibleCounts[g]} $\leftarrow 0$.

  \begin{columns}[T]
    \begin{column}{0.5\textwidth}
      \begin{itemize}
        \item \small Ràng buộc hay bị vi phạm $\to$ tăng phạt, \hl{kéo về vùng khả thi}
        \item \small Ràng buộc luôn thoả $\to$ giảm phạt, \hl{cho phép đi qua vùng không khả thi} để tìm nghiệm tốt hơn
      \end{itemize}
    \end{column}
    \begin{column}{0.5\textwidth}
      \centering\small
      \begin{tabular}{ll}
        \toprule
        Tham số & Giá trị \\
        \midrule
        $\lambda$ ban đầu & 1 \\
        $\lambda^{\min}$ / $\lambda^{\max}$ & 0.001 / 1000 \\
        $L_\lambda$ & 20 \\
        $\rho_{\min}$ / $\rho_{\max}$ & 0.20 / 0.80 \\
        $\gamma_\uparrow$ / $\gamma_\downarrow$ & 1.50 / 0.85 \\
        \bottomrule
      \end{tabular}
    \end{column}
  \end{columns}
\end{frame}

\begin{frame}{Ruin \& Recreate — Ruin}
  \small Ruin \& Recreate \hl{không phải move thông thường} trong candidate pool. Nó chỉ được gọi khi $h_{Div} \ge H_{Div}$.

  \vspace{2mm}
  \textbf{Số khách bị loại:}
  $$q_{\text{ruin}} = \max\{1, \lfloor \rho_{\text{ruin}}\, n \rfloor\}, \qquad \rho_{\text{ruin}} = 0.15$$

  \textbf{Chọn hỗn hợp:}
  \begin{itemize}
    \item \small Khoảng $1/3$: khách có đóng góp vi phạm ($V_{TW} + V_W$) lớn nhất
    \item \small Khoảng $1/3$: khách trên phương tiện critical (tăng lên $q/2$ nếu nghiệm đang khả thi)
    \item \small Phần còn lại: chọn ngẫu nhiên
  \end{itemize}

  \vspace{1mm}
  \small Gỡ các khách này khỏi nghiệm, xoá trip rỗng, đánh giá lại nghiệm bộ phận.
\end{frame}

\begin{frame}[fragile]{Ruin \& Recreate — Recreate}
  \begin{columns}[T]
    \begin{column}{0.48\textwidth}
      \textbf{Các khách bị loại được:}
      \begin{enumerate}
        \item \small Sắp theo $l_i$ tăng dần
        \item \small Nhiễu ngẫu nhiên nhẹ (đổi chỗ 2 khách liền kề với xác suất 0.1)
        \item \small Chèn lại bằng \hl{best insertion} trên toàn hệ thống
      \end{enumerate}
      \textbf{Mỗi khách được thử tại:}
      \begin{itemize}
        \item \small mọi vị trí trong mọi trip tương thích
        \item \small trip mới tại mọi vị trí trong chuỗi trip
      \end{itemize}
      \small Không ép khách vào \texttt{trucks[0]}.
    \end{column}
    \begin{column}{0.52\textwidth}
      \begin{block}{RUIN\_RECREATE (rút gọn)}
        \begin{pseudo}
for attempt = 1 .. 5:
\ind partial $\leftarrow$ s; gỡ q khách
\ind for each khách i (đã sắp):
\ind\ind cand $\leftarrow$ EVALUATE\_MOVE(mọi vị trí)
\ind\ind if cand rỗng: thất bại, break
\ind\ind partial $\leftarrow$ SELECT\_BEST(cand)
\ind if thành công: return partial
\ind $\rho$ $\leftarrow$ max(0.05, 0.8 $\rho$)
return FAILURE
        \end{pseudo}
      \end{block}
      \small Thành công: $s \leftarrow$ nghiệm mới, \hl{xoá danh sách tabu}, $h_{Div} \leftarrow 0$.
    \end{column}
  \end{columns}
\end{frame}

% ===================================================================
\section{Thuật toán tổng thể}

\begin{frame}[fragile]{ADAPTIVE\_TABU\_SEARCH (1/2) — Khởi tạo và đa dạng hoá}
  \begin{columns}[T]
    \begin{column}{0.5\textwidth}
      \begin{pseudo}
1. s $\leftarrow$ s0 (Greedy Insertion)
2. H $\leftarrow$ max(1, C(s0), max l\_i)
3. $\lambda_Q = \lambda_D = \lambda_{TW} = \lambda_W$ $\leftarrow$ 1
4. EVALUATE\_SOLUTION(s, $\lambda$)
5. if IS\_FEASIBLE(s): s* $\leftarrow$ s
\ind else: s$^I$ $\leftarrow$ s
6. tabuUntil $\leftarrow$ $\emptyset$
7. feasibleCounts $\leftarrow$ 0
8. iter, hDiv, hStop $\leftarrow$ 0
      \end{pseudo}
    \end{column}
    \begin{column}{0.5\textwidth}
      \begin{pseudo}
9. while iter < Nmax and time < Tlim
\ind\ind and hStop < HStop:
10.\ iter $\leftarrow$ iter + 1
11.\ if hDiv $\ge$ HDiv:
12.\ \ind r $\leftarrow$ RUIN\_RECREATE(s, ...)
13.\ \ind if r $\ne$ FAILURE:
14.\ \ind\ind s $\leftarrow$ r; clear tabuUntil
\ind\ind\ind hDiv $\leftarrow$ 0
15.\ \ind UPDATE\_BEST\_SOLUTIONS
16.\ \ind cập nhật hStop
17.\ \ind continue
      \end{pseudo}
    \end{column}
  \end{columns}
\end{frame}

\begin{frame}[fragile]{ADAPTIVE\_TABU\_SEARCH (2/2) — Vòng lặp lân cận}
  \begin{columns}[T]
    \begin{column}{0.5\textwidth}
      \begin{pseudo}
18. pool $\leftarrow$ BUILD\_CANDIDATE\_POOL(s, s*, $\lambda$)
19-23. admissible $\leftarrow$ \{c : not IS\_TABU(c)
\ind\ind\ind or SATISFIES\_ASPIRATION(c)\}
24. if admissible = $\emptyset$:
25. \ind mở rộng candidate list, sinh lại
26. \ind vẫn tabu hết: giải phóng thuộc tính
\ind\ind sắp hết hạn sớm nhất, lọc lại
27. \ind không có move hợp lệ:
\ind\ind hDiv $\leftarrow$ HDiv; continue
      \end{pseudo}
    \end{column}
    \begin{column}{0.5\textwidth}
      \begin{pseudo}
28. c $\leftarrow$ SELECT\_BEST\_CANDIDATE(admissible)
29. s $\leftarrow$ c.solution
30. REGISTER\_TABU(c, iter, $\tau_0$)
31. improved $\leftarrow$ UPDATE\_BEST\_SOLUTIONS
32. cập nhật hDiv, hStop
33-36. cập nhật feasibleCounts
37. if iter mod L$_\lambda$ = 0:
38. \ind UPDATE\_PENALTIES($\lambda$)
39. return s* nếu có
40. else: báo "No feasible solution"
\ind\ind return s$^I$ (chẩn đoán)
      \end{pseudo}
    \end{column}
  \end{columns}
\end{frame}

\begin{frame}{Điều kiện dừng và các chi tiết bảo vệ}
  \textbf{Dừng} khi một trong ba điều kiện xảy ra:
  \begin{itemize}
    \item \small đủ $N_{\max}$ vòng lặp
    \item \small hết thời gian $T_{\lim}$
    \item \small $h_{Stop} \ge H_{Stop}$: quá lâu không cải thiện record
  \end{itemize}

  \vspace{2mm}
  \textbf{Tránh tìm kiếm bị chặn hoàn toàn bởi tabu (bước 24–27)}
  \begin{itemize}
    \item \small Mở rộng nguồn sinh move ra \hl{mọi khách, mọi trip} rồi sinh lại pool.
    \item \small Vẫn tabu hết: giải phóng dần các thuộc tính sắp hết hạn sớm nhất và lọc lại.
    \item \small Không còn move hợp lệ nào: ép đa dạng hoá ở vòng sau ($h_{Div} \leftarrow H_{Div}$).
  \end{itemize}

  \vspace{2mm}
  \begin{alertblock}{Bổ sung khi cài đặt}
    \small $h_{Div}$ được đặt lại về 0 sau \emph{mỗi} lần thử Ruin \& Recreate, kể cả khi thất bại. Nếu không, lần thử thất bại sẽ bị gọi lại ở mọi vòng và Tabu Search không bao giờ chạy tiếp.
  \end{alertblock}
\end{frame}

% ===================================================================
\section{Cài đặt và thử nghiệm}

\begin{frame}{Tham số thuật toán}
  \centering
  \begin{tabular}{llr}
    \toprule
    Ký hiệu & Ý nghĩa & Giá trị \\
    \midrule
    $N_{\max}$ & số vòng lặp tối đa & 2000 \\
    $T_{\lim}$ & giới hạn thời gian & tuỳ thí nghiệm \\
    $H_{Stop}$ & số vòng không cải thiện thì dừng & 500 \\
    $H_{Div}$ & số vòng không cải thiện thì Ruin \& Recreate & 30 \\
    $\tau_0$ & tabu tenure cơ sở, $\tau \sim U[7, 14]$ & 7 \\
    $L_\lambda$ & chu kỳ cập nhật $\lambda$ & 20 \\
    $\rho_{\text{ruin}}$ & tỷ lệ khách bị loại khi Ruin & 0.15 \\
    -- & số lần thử Ruin \& Recreate & 5 \\
    -- & ngưỡng vi phạm tối đa của candidate & 5.0 \\
    -- & tỷ lệ nguồn ngẫu nhiên khi chọn thành phần & 30\% \\
    \bottomrule
  \end{tabular}

  \vspace{3mm}
  \small Giá trị hiện tại là giá trị đặt tạm; sẽ được hiệu chỉnh bằng thực nghiệm.
\end{frame}

\begin{frame}{Cấu trúc mã nguồn (C++17)}
  \begin{columns}[T]
    \begin{column}{0.5\textwidth}
      \footnotesize
      \begin{tabular}{@{}ll@{}}
        \toprule
        File & Nội dung \\
        \midrule
        \texttt{instance.hpp} & đọc dữ liệu, quy đổi $L_D$ \\
        \texttt{solution.hpp} & Trip, Vehicle, Solution \\
        \texttt{schedule.hpp} & tương thích tĩnh, tính lịch \\
        \texttt{evaluate.hpp} & $H$, vi phạm, $F_\lambda$ \\
        \texttt{construction.hpp} & nghiệm khởi tạo \\
        \texttt{operators.hpp} & 6 toán tử \\
        \texttt{evaluate\_move.hpp} & đánh giá move \\
        \bottomrule
      \end{tabular}
    \end{column}
    \begin{column}{0.5\textwidth}
      \footnotesize
      \begin{tabular}{@{}ll@{}}
        \toprule
        File & Nội dung \\
        \midrule
        \texttt{select\_components.hpp} & chọn thành phần \\
        \texttt{candidate\_pool.hpp} & candidate pool \\
        \texttt{tabu.hpp} & tabu, aspiration \\
        \texttt{select\_move.hpp} & chọn move \\
        \texttt{strategic\_oscillation.hpp} & cập nhật $\lambda$ \\
        \texttt{ruin\_recreate.hpp} & Ruin \& Recreate \\
        \texttt{tabu\_search.hpp} & vòng lặp chính \\
        \texttt{validator.hpp} & kiểm tra nghiệm \\
        \bottomrule
      \end{tabular}
    \end{column}
  \end{columns}
\end{frame}

\begin{frame}{Kiểm chứng và tối ưu cài đặt}
  \textbf{Kiểm tra nghiệm cuối (validator độc lập)}
  \begin{itemize}
    \item \small Tự mô phỏng lại toàn bộ lịch từ toạ độ và dữ liệu thô, không dùng lại hàm của thuật toán.
    \item \small Kiểm tra: mỗi khách đúng 1 lần, $C_1$ không đi drone, tải trọng, quãng đường bay $\le L_D$, $a_i \le l_i$, thời gian chờ $\le L_w$.
    \item \small Đối chiếu số liệu đã lưu (makespan, quãng đường, thời điểm) với giá trị mô phỏng lại.
  \end{itemize}

  \vspace{2mm}
  \textbf{Incremental evaluation}
  \begin{itemize}
    \item \small Mỗi xe lưu sẵn vi phạm thô của riêng nó $\to$ tổng hợp nghiệm chỉ cần cộng theo số xe.
    \item \small Chỉ tính lại lịch các xe bị move chạm vào.
    \item \small Copy-on-write: sao chép nghiệm chỉ sao chép con trỏ xe; xe nào bị sửa mới được nhân bản.
  \end{itemize}
\end{frame}

\begin{frame}{Kết quả chạy thử}
  \centering
  \begin{tabular}{lccccc}
    \toprule
    Instance & $n$ & $K$ / $D$ & Số vòng & $C_{\max}$ & Validator \\
    \midrule
    \texttt{6\_5\_1}     & 6  & 1 / 1 & 502 & 699.67  & hợp lệ \\
    \texttt{easy\_test}  & 8  & 2 / 2 & 500 & 43.87   & hợp lệ \\
    \texttt{10\_20\_4}   & 10 & 1 / 1 & 504 & 6469.04 & hợp lệ \\
    \bottomrule
  \end{tabular}

  \vspace{3mm}
  \begin{itemize}
    \item \small Cả ba instance đều tìm được nghiệm khả thi; validator độc lập xác nhận nghiệm thoả mọi ràng buộc.
    \item \small Thời gian chạy dưới 1 giây; thuật toán dừng do $h_{Stop} \ge H_{Stop} = 500$.
    \item \small Instance lớn \texttt{100\_40\_4} ($n = 100$, 4 truck, 4 drone): đang tiếp tục thử nghiệm và hiệu chỉnh tham số.
  \end{itemize}
\end{frame}

% ===================================================================
\section{Kết luận}

\begin{frame}{Kết luận và hướng phát triển}
  \textbf{Đã thực hiện}
  \begin{itemize}
    \item \small Mô hình hoá bài toán truck--drone nhiều chuyến, có time window và thời gian chờ hàng.
    \item \small Cài đặt đầy đủ Adaptive Tabu Search: nghiệm khởi tạo, 6 toán tử, 3 loại thuộc tính tabu, aspiration, Strategic Oscillation, Ruin \& Recreate.
    \item \small Validator độc lập và incremental evaluation.
  \end{itemize}

  \vspace{2mm}
  \textbf{Hướng phát triển}
  \begin{itemize}
    \item \small Hiệu chỉnh tham số ($H_{Div}$, $\tau_0$, $\rho_{\text{ruin}}$, $T_{\lim}$) bằng thực nghiệm.
    \item \small Thử nghiệm trên bộ dữ liệu lớn, so sánh với các phương pháp khác.
    \item \small Chuẩn hoá bộ dữ liệu test theo đúng định nghĩa $L_D$ là quãng đường.
  \end{itemize}
\end{frame}

\begin{frame}
  \centering
  \Huge Cảm ơn thầy cô đã lắng nghe!
\end{frame}

\end{document}
